% Insert once after odd-q-packet-template.tex and before the local modules.
\subsection{Finite rows with a two-edge root arm}
\label{sec:odd-common:w2}

\begin{lemma}[Continuation from an anchored first row]
\label{lem:odd-common:anchored-continuation}
Let \(s\ge1\), and suppose an actual terminal has the positive upper bound
\(TF\le e^\alpha[\bar f_s+\mu BR+D]\), where \(R\) is real.
Put \(D_i=\E B^iR\). For every integer \(L\ge1\),
\begin{equation}\label{eq:odd-common:anchored-continuation}
 0<e^{-\alpha}T^LF\le
 \bar f_{s+L-1}+\mu B^LR+D\bar f_{L-1}
       +\mu\sum_{i=1}^{L-1}D_i\bar f_{L-1-i}.
\end{equation}
The constants \(\alpha,\mu,D\) are those of the proved first row.
\end{lemma}
\begin{proof}
Apply \(T^{L-1}\) to that row and use
\eqref{eq:oc:finite-duhamel} with the vector \(BR\).

\end{proof}

For the actual two-edge arm \(W_2=1-u-Bu\), any positive upper bound
\(T^LF\le e^{-\Gamma}[Cq+BV+c]\) gives
\begin{equation}\label{eq:odd-common:w2-majorant}
 \log(W_2T^LF)\le-\Gamma+C+c-1
                   +B[V-u+(C-1)\one].
\end{equation}
This follows by adding \(\log W_2\le-u-Bu\) to
\(\log(Cq+BV+c)\le Cq+BV+c-1\).
The complete bracket, rather than each summand, is required to be positive.

\begin{lemma}[The ordinary terminal at a \(W_2\) root]
\label{lem:odd-common:w2}
Suppose \(F>0\), \(F\le e^{R-\Gamma}\), and the actual remainder
satisfies \(e^R\le1+R+(c_T/2)R^2\).
Write \(\sigma,D_i,\delta,\epsilon\) as in
\eqref{eq:oc:fourth-data}, with \(c=c_T\).
Then \(\log(W_2T^LF)\le a+B(V_0+c_0\one)\), where
\begin{equation}\label{eq:odd-common:w2-table}
\begin{array}{c|c|c|c}
 L&a&V_0&c_0\\ \hline
 1&-\Gamma+\delta&R-u&0\\
 2&-\Gamma+\epsilon&BR&\delta\\
 3&-\Gamma+\epsilon-d+D_2&(B+\delta)u+B^2R&\epsilon\\
 L\ge4&-\Gamma+2\delta_++2(D_1)_+-2hd
          &(G_{L-1}-1)u+B^{L-1}R&-hd
\end{array}
\end{equation}
The last row assumes \(h\ge0\) and
\(\gamma_j=\beta_j-D_{j-1}\ge hd\) for every \(j\ge3\).
With \(n=L-1\), its two means are
\begin{equation}\label{eq:odd-common:w2-long-means}
 \E BV_0=d-\gamma_{L+1}\le(1-h)d,
 \qquad \E V_0=-\gamma_{n+1}\le-hd.
\end{equation}
At a unit-length row whose actual arm is \(f_2\), replace only
\(R-u\) by \(R+u\).
If the actual root-pair product has envelope
\(P\le e^{-\Gamma_r+BV_r}\), including its full constant,
\begin{equation}\label{eq:odd-common:w2-complete}
\begin{gathered}
 V=V_0+c_0\one+V_r,\qquad S=a-\Gamma_r+\E BV,\\
 \log\E[(PW_2T^LF)^4]\le4S+\tfrac{17}{4}\|V\|^2.
\end{gathered}
\end{equation}
For the unit \(f_2\) exception the same conclusion uses that actual arm.
\end{lemma}
\begin{proof}
The finite positive majorant is \eqref{eq:oc:fourth-duhamel}; its
derivation needs only the displayed remainder bound and
\(T(R^2)\le2\sigma\).
Expand its first three lengths and add \(\log W_2\le-u-Bu\).
For \(L\ge4\), use \eqref{eq:odd-common:long-parameters}; adding
\(\log W_2\) subtracts one \(u\) from its primitive and gives
the last row of \eqref{eq:odd-common:w2-table}.
The signed means follow from \eqref{eq:odd-common:coordinates}.
The unit exception uses \(\log f_2\le-u+Bu\) instead.
Finally add the actual pair envelope and apply
Lemma~\ref{lem:odd-common:fourth} to the whole vector \(V\).
There is no independent unit-arm curvature payment at this root.
\end{proof}
